\subsection{Entwicklungsbetrieb}
\label{subsec:entwicklungsbetrieb}

Abbildung~\ref{fig:developer-workflow} zeigt den profilabhängigen Weg von
Diagnose und Installation bis zu Qualitätssicherung und Auslieferung.

\begin{figure}[htbp]
  \centering
  \resizebox{0.72\textwidth}{!}{%
  \begin{tikzpicture}[node distance=5mm and 11mm]
    \node[casePrimary,minimum width=50mm] (doctor) {\texttt{doctor}};
    \node[caseBox,minimum width=50mm,below=of doctor] (install) {\texttt{install}};
    \node[caseBox,minimum width=50mm,below=of install] (run) {\texttt{run}};
    \node[caseBox,minimum width=50mm,below=of run] (develop) {Produktcode bearbeiten};
    \node[caseBox,minimum width=50mm,below=of develop] (quality) {\texttt{quality}};
    \node[caseBox,minimum width=50mm,below=of quality] (test) {\texttt{test --suite all}};
    \node[caseBox,minimum width=50mm,below=of test] (build) {passendes \texttt{build}-Ziel};
    \node[caseOptional,minimum width=50mm,below=of build] (release) {\texttt{release check}};

    \node[caseOptional,minimum width=38mm,right=of run] (desktop)
      {Desktopprofil\\\texttt{tauri run}};
    \node[caseOptional,minimum width=38mm,right=of build] (targets)
      {Web / Desktop /\\Container nach Profil};

    \draw[caseFlow] (doctor) -- (install);
    \draw[caseFlow] (install) -- (run);
    \draw[caseFlow] (run) -- (develop);
    \draw[caseFlow] (develop) -- (quality);
    \draw[caseFlow] (quality) -- (test);
    \draw[caseFlow] (test) -- (build);
    \draw[caseFlow] (build) -- (release);
    \draw[caseSecondaryFlow] (quality.west) to[out=180,in=180,looseness=1.5] (develop.west);
    \draw[caseSecondaryFlow] (run) -- (desktop);
    \draw[caseSecondaryFlow] (build) -- (targets);
  \end{tikzpicture}}
  \caption{Grundstruktur des Entwicklungsworkflows}
  \label{fig:developer-workflow}
  \smallskip
  {\footnotesize Quelle: Eigene Darstellung auf Grundlage von
  \textcite{kleiveist2026tooling}, \textcite{kleiveist2026release} und
  \textcite{kleiveist2026quality-policy}.\par}
\end{figure}


Nach einer Änderung liegt \texttt{quality} vor Tests und Build. Warnungen
bleiben im Report sichtbar und führen zurück in den Bearbeitungszyklus, ohne
allein zu blockieren; ein aktiver Fehler beendet den Befehl mit Exitcode~1.
Dieser lokale Einstieg entspricht dem Quality-Aufruf in Core-, Profil- und
Release-Workflows, nicht jedoch einer Behauptung identischer
Ausführungsumgebungen.

\texttt{run} startet Vite und nur bei aktivem Backend zusätzlich Uvicorn. Im
Vordergrund beendet ein Abbruch die erfassten Prozessgruppen; im abgekoppelten
Modus werden PID- und Logdaten gespeichert. \texttt{stop} beendet erfasste
Prozesse und standardmäßig auch Listener auf den konfigurierten Ports;
\texttt{--tracked-only} beschränkt dies auf den eigenen Zustand
\parencite{kleiveist2026control,kleiveist2026tooling}.
Im Vordergrund werden die Ausgaben aktiver Dienste zusammengeführt. Der
abgekoppelte Modus hält dagegen Prozess- und Logdaten unter
\path{tools/.runtime}; diese Zustandsablage bildet die Grundlage für
\texttt{stop}.

Der getrennte Desktopweg delegiert an \texttt{tauri dev} und entfernt bekannte
serverseitige Geheimnisse aus der Kindprozessumgebung. Er unterstützt Vorder-
und Hintergrundbetrieb. Die Konsole führt lediglich durch dieselben
CLI-Workflows.
Damit bleibt der native Entwicklungsweg vom allgemeinen Vite-/Uvicorn-Betrieb
getrennt, ohne eine zweite Befehlslogik einzuführen.
\beginnerinput{workspace/statement/500_tooling_workflow/540}
